% Beamer-Klasse mit 4:3-Seitenverhältnis
%\documentclass{beamer}
% Beamer-Klasse mit 16:9-Seitenverhältnis
% Anmerkung: Bei der Angabe 32:18 bleibt die Höhe der Folien konstant.
\documentclass[aspectratio=3218]{beamer}

% Pakete laden
\input{pakete}

% Stil des Lehrstuhls für Elektromagnetische Verträglichkeit laden
%\usepackage{beamer_emv}
% Stil des Instituts für Medizintechnik laden
%\usepackage{beamer_imt}
% Stil der Fakultät für Elektro- und Informationstechnik laden
%\usepackage{beamer_feit}
% Stil der Otto-von-Guericke-Universität laden
%\usepackage{beamer_ovgu}
\usepackage{beamer_ovgu_169}
% Stil des Stimulate e.V. laden
%\usepackage{beamer_stimulate}

\title[Kurztitel]{Titel der Präsentation}
\author{Max Mustermann}
\institute[Lehrstuhl EMV]{
	Lehrstuhl für Elektromagnetische Verträglichkeit \\
	Institut für Medizintechnik \\
	Otto-von-Guericke-Universität, Magdeburg
}
\mode<presentation>{\keywords{Schlüsselwörter durch Komma getrennt}}
\date[01.01.2023]{Datum der Präsentation, \zB 1. Januar 2023}

\begin{document}

\begin{frame}
	\maketitle
	% \maketitle funktioniert auch im Article-Modus
	% \titlepage funktioniert nur bei Präsentationen
\end{frame}

\begin{frame}[label=inhalt]{Gliederung}
	\tableofcontents
\end{frame}

\section{Einleitung}

% nur für die Miniframe-Navigation
\subsection*{}

\begin{frame}
	\frametitle<presentation>{Einleitung}
	\begin{block}{Sinn und Zweck einer Abschlusspräsentation:}
		\begin{itemize}
			\item Darstellung der wichtigsten Ergebnisse der Arbeit
			\item Interesse für die Arbeit/den Bericht zu wecken
			\item Ziel: lehrreich und unterhaltsam zugleich
		\end{itemize}
	\end{block}
	\begin{block}{Vorteile von \LaTeX\ mit der \textsc{beamer}-Klasse:}
		\begin{itemize}
			\item sehr einfach, wenn die Arbeit bereits in \LaTeX\ erstellt wurde
			\item verschiedene Darstellungsvarianten zur Auswahl
			\item Handzettel können einfach mit \emph{beamerarticle} erzeugt werden
		\end{itemize}
	\end{block}
\end{frame}

\section{Hauptteil}

\subsection{Gleichungen}

\begin{frame}
	\frametitle<presentation>{Gleichungen}
	\begin{block}{Satz des Pythagoras:}
		\begin{equation}
			a^2 + b^2 = c^2 \label{eq:pythogoras}
		\end{equation}	\end{block}
	\begin{block}{Es folgt, dass:}
		\begin{align}
			a^2 &= c^2 - b^2 \\
			b^2 &= c^2 - a^2
		\end{align} 	
 	\end{block}
\end{frame}

\subsection{Abbildungen}

\begin{frame}
	\frametitle<presentation>{Abbildungen}
	\begin{figure}
		\centering
			\includegraphics[height=5cm]{bilder/unilogo}
		\caption{Altes Logo der Universität}
		\label{fig:unilogo}
	\end{figure}
\end{frame}

\begin{frame}{Diagramme}
	\begin{figure}
		\centering
			\begin{tikzpicture}
				\begin{axis}[
					beamer,
					xlabel={Zeit, $t$ (in \si{\milli\second})},
					ylabel={Spannung, $u(t)$ (in \si{\volt})},
					xmin=0,xmax=20,
					ymin=-350,ymax=350,
					legend pos=south west,
					% Beispiel für das Einbinden einer Linie durch Angabe einer Formel
					]
					\addplot+[
						domain=0:20,
						samples=101,
					] {sin(deg(x*2*pi/20))*sqrt(2)*230};
					\addlegendentry{sinusförmig};
					\addplot+[
						domain=0:20,
						samples=101,
					] {cos(deg(x*2*pi/20))*sqrt(2)*230};
					\addlegendentry{cosinusförmig};
				\end{axis}
			\end{tikzpicture}
		\caption{Harmonischer Zeitverlauf einer Spannung mit einer Frequenz von \SI{50}{\hertz} und einem Effektivwert von \SI{230}{\volt}}
		\label{fig:u_t_sinus}
	\end{figure}
\end{frame}

\subsection{Tabellen}

\begin{frame}
	\frametitle<presentation>{Tabellen}
	\begin{table}
		\caption{Beispieltabelle}
		\label{tab:beispiel}
		\centering
			\begin{tabular}{cc}
				\toprule
				Variable & Bedeutung \\
				\midrule
				$t$ & Zeit \\
				$U$ & Spannung \\
				\bottomrule
			\end{tabular}
	\end{table}
\end{frame}

\begin{frame}{Zitate}
	\begin{block}{Keine Kurzzitate benutzen:}
		\begin{itemize}
			\item keine Zitate wie~[1] einfügen
			\item niemand erinnert sich an die Nummer, wenn das Literaturverzeichnis gezeigt wird
			\item stattdessen Vollzitate
		\end{itemize}
	\end{block}
	\begin{block}{Beispiel für ein Vollzitat:}
		\small\fullcite{hering:berichte}
	\end{block}
\end{frame}

\section{Zusammenfassung}

% nur für die Miniframe-Navigation
\subsection*{}

\begin{frame}
	\frametitle<presentation>{Zusammenfassung}
	\begin{block}{Ergebnisse:}
	  \begin{itemize}
	  	\item Kurzfassung der wichtigsten Ergebnisse
	  	\item Aufzeigen nicht gelöster Probleme und offener Fragen
	  	\item eventuell Ausblick auf weitere geplante Forschungsarbeiten
		\end{itemize}
	\end{block}
	\begin{block}{Fragen der Zuhörer:}
	 \begin{itemize}
	 	\item genügend Zeit für Fragen einplanen
	 	\item optional: zusätzliche Folien für wahrscheinliche Fragen
	 \end{itemize}
	\end{block}	
\end{frame}

% nur für die Miniframe-Navigation
\section*{}

\begin{frame}<beamer>{}
	\begin{center}
		Vielen Dank für Ihre Aufmerksamkeit!
	\end{center}
	\begin{center}
		Gibt es Fragen?
	\end{center}
\end{frame}

\end{document}
